%=====================================================================
% PROJETO PREDIAL - ENTRADA, MEDICAO E INTERFACE COM A DISTRIBUIDORA
%=====================================================================

\section{Entrada, medição e padrão Equatorial Maranhão}

\begin{frame}{A fronteira da distribuidora precisa estar explícita}
\centering
\begin{tikzpicture}[
  b/.style={draw=corPrim,fill=corDef!7,rounded corners,minimum width=2.05cm,minimum height=0.72cm,align=center,font=\scriptsize},
  a/.style={-{Stealth},thick,corPrim}]
\node[b](rede)at(-5.0,0){Rede de\\distribuição};
\node[b](pc)at(-2.5,0){Ponto de\\conexão};
\node[b](ramal)at(0,0){Ramal de\\conexão};
\node[b](med)at(2.5,0){Medição e\\proteção geral};
\node[b](qg)at(5.0,0){Origem da instalação\\e QD/QGBT};
\draw[a](rede)--(pc);\draw[a](pc)--(ramal);\draw[a](ramal)--(med);\draw[a](med)--(qg);
\draw[corAccent,very thick,dashed] (1.25,-1.0)--(1.25,1.0);
\node[font=\tiny,corAccent,align=center] at (-1.0,-1.25){critérios da distribuidora};
\node[font=\tiny,corAccent,align=center] at (3.5,-1.25){instalação do consumidor};
\end{tikzpicture}
\vspace{5mm}
\begin{caixaObs}
Responsabilidade física, propriedade dos componentes e limite elétrico não devem ser presumidos. Confirmar a definição na norma vigente e no orçamento de conexão.
\end{caixaObs}
\end{frame}

\begin{frame}[shrink=4]{Documentos vigentes usados no Maranhão}
\scriptsize
\begin{tabularx}{\textwidth}{>{\bfseries}p{2.8cm}p{2.0cm}X}
\toprule
Documento & Revisão consultada & Aplicação \\
\midrule
NT.00001.EQTL & Rev. 09, 29/05/2025 & unidade individual em BT, carga instalada até \SI{75}{kW} \\
NT.00004.EQTL & Rev. 08, 31/12/2025 & empreendimentos de múltiplas unidades consumidoras \\
Anexo I da NT.00004 & Rev. 08, 21/05/2026 & cálculo de queda de tensão em EMUC \\
NT.00020.EQTL & versão vigente & micro e minigeração distribuída \\
NT.00030.EQTL & versão vigente & caixas de medição e proteção \\
NT.00042.EQTL & versão vigente & conexão de estação de recarga de veículo elétrico \\
\bottomrule
\end{tabularx}
\vspace{2mm}
\begin{caixaAt}
Revisões podem mudar durante o projeto. Registrar a data da consulta e conferir novamente antes de protocolar, comprar a caixa e executar o padrão.
\end{caixaAt}
\end{frame}

\begin{frame}[shrink=4]{Tabela 220/380 V: enquadramento inicial da unidade individual}
\scriptsize
\begin{tabularx}{\textwidth}{>{\bfseries}p{2.3cm}p{2.2cm}c c X}
\toprule
Carga instalada & Fornecimento & DG & Cu cliente & Observação \\
\midrule
até \SI{5}{kW} & monofásico 220 V & 25 A & \SI{4}{mm^2} & faixa inicial \\
5,1--\SI{8}{kW} & monofásico 220 V & 40 A & \SI{6}{mm^2} & -- \\
8,1--\SI{12}{kW} & monofásico 220 V & 63 A & \SI{10}{mm^2} & limite monofásico urbano no MA \\
até \SI{12}{kW} & trifásico 380/220 V & 20 A & \SI{6}{mm^2} & opção pode gerar custo de adaptação \\
12,1--\SI{24}{kW} & trifásico 380/220 V & 40 A & \SI{6}{mm^2} & -- \\
24,1--\SI{38}{kW} & trifásico 380/220 V & 63 A & \SI{10}{mm^2} & faixa do projeto residencial I \\
38,1--\SI{48}{kW} & trifásico 380/220 V & 80 A & \SI{16}{mm^2} & -- \\
48,1--\SI{60}{kW} & trifásico 380/220 V & 100 A & \SI{25}{mm^2} & -- \\
60,1--\SI{75}{kW} & trifásico 380/220 V & 125 A & \SI{35}{mm^2} & limite da tabela BT \\
\bottomrule
\end{tabularx}
\vspace{1mm}
\tiny Fonte técnica: Tabela 1 da NT.00001.EQTL Rev.09. Reproduzidos apenas os campos necessários ao enquadramento didático; consultar a tabela integral para ramal, eletrodutos, aterramento e notas.
\end{frame}

\begin{frame}{Exemplo: residência com \SI{33.96}{kW} instalados}
\small
\begin{columns}[T,onlytextwidth]
\column{0.48\textwidth}
\begin{align*}
P_{\mathrm{inst}}&=\SI{33.96}{kW}\\
24{,}1&<33{,}96\leq38
\end{align*}
\begin{caixaEx}[title=Enquadramento]
\begin{itemize}
\item fornecimento: 380/220 V, 3F+N;
\item disjuntor geral: 63 A;
\item condutor Cu do cliente: mínimo \SI{10}{mm^2};
\item verificar demais colunas e notas da tabela.
\end{itemize}
\end{caixaEx}
\column{0.50\textwidth}
\begin{caixaAt}[title=Não encerrar o projeto aqui]
A tabela da distribuidora dimensiona o padrão sob hipóteses próprias. O alimentador interno ainda deve ser recalculado por:
\begin{itemize}
\item capacidade de condução;
\item queda de tensão;
\item curto-circuito e proteção;
\item neutro, PE e método real.
\end{itemize}
\end{caixaAt}
\end{columns}
\end{frame}

\begin{frame}{Padrão aéreo: sequência física dos componentes}
\centering
\begin{tikzpicture}[scale=0.88,every node/.style={font=\scriptsize}]
\draw[very thick] (-5.2,-1.7)--(-5.2,3.3);
\draw[very thick] (-5.6,2.8)--(-4.8,2.8);
\draw[corPrim,thick] (-5.2,2.8).. controls (-2.5,2.5) and (-1.6,2.2)..(0,2.0);
\draw[very thick] (0,-1.7)--(0,3.0);
\node[draw=corPrim,fill=corDef!7,minimum width=1.8cm,minimum height=1.2cm,align=center](cx)at(0.95,0.75){caixa de\\medição};
\node[draw=corPrim,fill=corDef!7,minimum width=1.8cm,minimum height=0.8cm,align=center](dg)at(0.95,-0.45){proteção\\geral};
\draw[corPrim,thick] (0,2.0)--(cx.north);
\draw[corPrim,thick] (cx)--(dg);
\draw[corPrim,thick,-{Stealth}] (dg.east)--(3.15,-0.45)node[right]{alimentador};
\draw[corNorma,thick] (cx.south west)--(-0.65,-0.7)--(-0.65,-1.35)node[ground]{};
\node[align=center] at (-5.2,-2.1){poste da rede};
\node[align=center] at (0,-2.1){poste/muro do padrão};
\node[corAccent,align=center] at (-2.5,3.0){ramal de conexão};
\end{tikzpicture}
\begin{caixaObs}
Alturas, afastamentos, tipo de poste, eletroduto, caixa e ramal devem seguir os desenhos construtivos vigentes da Equatorial e a situação real do imóvel.
\end{caixaObs}
\end{frame}

\begin{frame}{Medição e proteção: acessibilidade sem acesso indevido}
\small
\begin{columns}[T,onlytextwidth]
\column{0.48\textwidth}
\begin{tikzpicture}[scale=0.83]
\draw[very thick,rounded corners] (0,0) rectangle (4.5,5.1);
\draw[thick] (0,2.0)--(4.5,2.0);
\node[draw=corPrim,circle,minimum size=1.35cm] at (2.25,3.55){kWh};
\node[draw=corPrim,fill=corDef!7,minimum width=1.7cm,minimum height=0.7cm] at (2.25,1.0){DG};
\node[font=\scriptsize] at (2.25,4.65){compartimento de medição};
\node[font=\scriptsize] at (2.25,1.65){proteção};
\draw[corAccent,thick,dashed] (0.25,2.25) rectangle (4.25,4.85);
\node[font=\tiny,corAccent] at (3.65,2.45){lacrável};
\end{tikzpicture}
\column{0.50\textwidth}
\begin{itemize}
\item caixa homologada para o tipo de fornecimento;
\item leitura e operação conforme acesso exigido;
\item compartimentos e pontos de lacre preservados;
\item condutores identificados e sem emendas indevidas;
\item raio de curvatura e espaço para terminais;
\item DG compatível com a tabela e capacidade de curto;
\item vedação, drenagem e resistência ao ambiente.
\end{itemize}
\end{columns}
\end{frame}

\begin{frame}{Aterramento do padrão não cria um PE isolado}
\centering
\begin{tikzpicture}[
  b/.style={draw=corPrim,fill=corDef!7,rounded corners,minimum width=2.25cm,minimum height=0.72cm,align=center,font=\scriptsize}]
\node[b](med)at(-4.2,1.2){Medição e DG};
\node[b](qd)at(-1.0,1.2){QD\\barras N e PE};
\node[b,draw=corNorma,fill=corNorma!7](bep)at(2.2,1.2){BEP\\equipotencialização};
\node[b,draw=corNorma,fill=corNorma!7](massas)at(5.2,1.2){Massas e partes\\condutivas};
\draw[-{Stealth},thick,corPrim](med)--(qd);
\draw[thick,corNorma](qd)--(bep)--(massas);
\draw[thick,corNorma](bep.south)--(2.2,-0.4)node[ground]{};
\node[font=\tiny,corNorma] at (3.0,-0.45){eletrodo(s)};
\end{tikzpicture}
\vspace{4mm}
\begin{caixaAt}
O condutor de proteção dos circuitos retorna ao barramento PE e ao BEP. Neutro e PE só se relacionam onde o esquema de aterramento determinar; não instalar pontes improvisadas nos quadros terminais.
\end{caixaAt}
\end{frame}

\begin{frame}{Ambiente costeiro altera materiais e manutenção}
\small
\begin{columns}[T,onlytextwidth]
\column{0.49\textwidth}
\begin{caixaNorma}[title=Até e após \SI{2}{km} da orla]
A NT.00001 diferencia materiais do ramal conforme a distância da orla marítima. O projeto deve registrar a localização e consultar a coluna correta da tabela.
\end{caixaNorma}
\column{0.49\textwidth}
\begin{itemize}
\item selecionar caixa, ferragens e fixações adequadas;
\item evitar pares galvânicos incompatíveis;
\item proteger entradas contra água e condensação;
\item aplicar torque e preparação conforme fabricante;
\item prever inspeção de corrosão e reaperto técnico.
\end{itemize}
\end{columns}
\begin{caixaAt}
Pintura ou fita não corrige material inadequado. A solução começa na especificação e segue na montagem e manutenção.
\end{caixaAt}
\end{frame}

\begin{frame}{EMUC: a unidade não define sozinha o empreendimento}
\centering
\begin{tikzpicture}[
  b/.style={draw=corPrim,fill=corDef!7,rounded corners,minimum width=2.0cm,minimum height=0.7cm,align=center,font=\scriptsize},
  a/.style={-{Stealth},thick,corPrim}]
\node[b](rede)at(0,3.4){Rede};
\node[b](med)at(0,2.1){Centro de medição};
\node[b](pr)at(0,0.8){Prumada / distribuição};
\node[b](u1)at(-3.6,-0.8){Unidade 1};
\node[b](u2)at(-1.2,-0.8){Unidade 2};
\node[b](u3)at(1.2,-0.8){Unidade 3};
\node[b](uc)at(3.6,-0.8){Área comum};
\draw[a](rede)--(med);\draw[a](med)--(pr);
\foreach \x in {u1,u2,u3,uc}{\draw[a](pr)--(\x);}
\end{tikzpicture}
\vspace{2mm}
\begin{caixaNorma}
Demanda do conjunto, centro de medição, prumadas, queda, proteção geral, espaço técnico e transferência de ativos seguem a NT.00004 e seus anexos vigentes.
\end{caixaNorma}
\end{frame}

\begin{frame}{Vistoria: conferir antes de solicitar a ligação}
\small
\begin{enumerate}
\item tipo de fornecimento, caixa e DG conferem com a carga declarada;
\item poste, ramal, eletrodutos, curvas e fixações seguem o desenho aplicável;
\item alturas, afastamentos e acesso ao medidor estão corretos;
\item condutores, terminais, identificação e raios de curvatura estão adequados;
\item aterramento, eletrodo e conexão estão instalados e acessíveis;
\item compartimentos lacráveis não foram invadidos;
\item padrão está íntegro, vedado, limpo e sem materiais improvisados;
\item dados do imóvel, carga e titular estão coerentes com a solicitação.
\end{enumerate}
\begin{caixaAt}
Não energizar a instalação interna como teste improvisado. Ensaios e energização seguem procedimento, instrumentos adequados e responsabilidades definidas.
\end{caixaAt}
\end{frame}

\begin{frame}{Erros que causam retrabalho no padrão}
\scriptsize
\begin{tabularx}{\textwidth}{>{\bfseries}p{3.4cm}X}
\toprule
Erro & Correção de projeto/obra \\
\midrule
Caixa comprada antes do enquadramento & fechar carga e tipo de fornecimento primeiro \\
DG fora da faixa normativa & usar tabela vigente e justificar caso especial \\
Poste ou eletroduto incompatível & seguir desenho construtivo aplicável \\
Medidor sem acesso regulamentar & reposicionar o padrão antes da execução final \\
Emenda em compartimento inadequado & prever comprimento e terminal correto \\
Neutro e PE unidos indevidamente & corrigir barras e esquema de aterramento \\
Material inadequado à orla & especificar conforme ambiente e tabela \\
Projeto interno divergente da carga declarada & revisar quadro, memorial e solicitação \\
\bottomrule
\end{tabularx}
\end{frame}
